\section{Discussion: Regular Local Rotation and Further Optimization}
\label{sec:rotation-discussion}

This section considers a regular change of leader \emph{inside the same
region}, while retaining BLAME-based recovery for stalled or equivocating
leaders. It does not install a new regional protocol. We distinguish a
conservative extension that completes a fixed application boundary before
handoff, an optional certificate-based opening optimization, and a replacement
of the regional engine by a protocol whose normal pipeline spans leaders.
The last option is how the HotStuff-family implementations examined below
organize ordinary leader progression. Their rules are evidence for design
choices, not proofs for the synchronous engine in this manuscript.

\subsection{Scope of the inherited quota-trigger experiment}
\label{sec:quota-experiment}
The copied Version 3 text of step V1 in \cref{alg:change} includes a tentative
condition that broadcasts BLAME after a predefined proposal count. We retain
that condition as an experiment, but the progress proof in
\cref{thm:progress} applies only to the original missing-progress and
equivocation triggers. In particular, its inference that an honest leader
is never blamed no longer holds under periodic quota-triggered BLAME.
The inherited safety argument relies on proof propagation, stopping old-view
votes, and safe inheritance; it does not require every exit to accuse a
faulty leader. Preserving those steps is therefore a plausible basis for
safety, but does not establish progress under repeated voluntary exits.

Proposal count alone also says nothing about completed timers. For example,
with $\delta_R=0.1$ ms and $\Delta_R=100$ ms, a healthy pipeline can propose
100 batches in approximately 20 ms, before its first $2\Delta_R$ commit
timer expires. An immediate quota-triggered quit can cancel every pending
timer in that view. This is a finite trace showing the missing progress
argument, not a claim of permanent deadlock for every quota interpretation.
If the condition is instead evaluated only at the existing
$(2p+4)\Delta_R$ deadlines, the proposal count may already exceed the
intended quota. A precise extension must specify both what is counted and
when the test is executed.

\subsection{How production HotStuff-family systems advance leaders}
\label{sec:production-rotation}
The sources below were inspected on 29 September 2026. Flow and Aptos code
citations identify repository commits; these are implementation snapshots,
not attestations of the binaries or feature flags used by every mainnet
validator. Official release descriptions are dated separately. In this
subsection $r$ denotes a comparison protocol's round, rather than this
manuscript's long-lived regional view $v$.

\paragraph{Flow: normal QC progression and TC-based recovery.}
Flow's HotStuff implementation describes itself as a Jolteon variant.
Votes for a block in view $r$ are sent to the leader of view $r+1$.
Its pacemaker has separate \texttt{ProcessQC} and \texttt{ProcessTC} paths.
The safety rules accept entry to $r$ directly when the parent QC has view
$r-1$; otherwise they require the preceding view's TC and check the QC
against the TC's newest QC. A local timeout alone does not advance the
view. Normal rotation consequently needs no timeout accusation, while
failed views have an explicit recovery certificate~\cite{flowcode}.
The finalization layer uses a certified child in the immediately following
view to finalize its parent. Progression and finalization are distinct
tests; a single parent QC is not interchangeable with the receipt of
\cref{sec:receipt}. Flow's engineering account explains that its active
pacemaker and shorter finalization rule were introduced together~\cite{flowjolteon}.

\paragraph{Aptos: the distinction is explicit in the round state.}
At commit \texttt{a19dae54b565}, \texttt{RoundState::process\_certificates}
sets the next round from the highest certificate and records its reason as
\texttt{QCReady} or \texttt{Timeout}. The inspected ordinary voting path
sends votes to the next round's elected proposer unless vote broadcasting
is configured. The same source also contains optimistic proposal handling:
when materializing an optimistic proposal, it checks both the parent QC's
round and the parent's identity. We therefore describe this as evidence
for certificate-driven progression, not as an exhaustive account of all
configured execution paths~\cite{aptosround,aptosmanager}.
The official 20 June 2025 Baby Raptr announcement describes an upgrade to
the Jolteon* and Quorum Store pipeline, including retrieval and verification
of batch contents before voting~\cite{aptosbaby}. An older DiemBFT account
alone is insufficient to characterize the current implementation. The
relevant lesson is that proposer selection, certified round advancement,
data readiness, and final commitment remain separate responsibilities.

\paragraph{Monad: next-leader collection also needs a recovery design.}
MonadBFT's documented normal path sends votes to the next scheduled leader,
which assembles the parent QC and proposes its child. Later certified
extensions finalize earlier blocks; speculative and full finality are
distinguished~\cite{monaddocs}. Its 21 November 2025 recovery update describes
backup votes sent to the current leader as well as the next leader, backup
QC dissemination, and votes carried in timeout messages that may reconstruct
a missing QC. More difficult cases retain reproposal and no-endorsement
certificates~\cite{monadrecovery}. These mechanisms address a specific risk:
a successor can withhold the votes it collects and attempt to discard its
predecessor's unfinalized tail. The current MonadBFT paper treats this
tail-forking resistance as a property beyond ordinary finalized-prefix
safety~\cite{monadpaper}. A transplant of next-leader vote routing without
the corresponding recovery rules would omit the main issue at the boundary.

\begin{table}[htbp]\centering\small
\begin{tabularx}{\linewidth}{@{}>{\raggedright\arraybackslash}p{.17\linewidth}Y Y@{}}\toprule
System & Ordinary advancement & Failure path and lesson\\\midrule
Flow & Parent QC enables the following view and its leader. & TC proves a skipped
view; voting rules constrain which certified parent may be extended.\\
Aptos snapshot & Explicit QC-ready versus timeout-driven round events. & Recovery
and optimistic proposal handling retain certificate checks and data prerequisites.\\
MonadBFT & Next leader collects votes and carries the resulting QC forward. &
Backup collection, timeout evidence and reproposal rules protect supported tails.\\
Current regional engine & QC enables another proposal by the \emph{same}
leader; timers establish commitment. & BLAME/equivocation exits the view;
STATUS and NEWVIEW recover a safe anchor.\\\bottomrule
\end{tabularx}
\caption{The transferable distinction is successful advancement versus failure
recovery. The commit rules, certificate meanings and fault assumptions are
specific to each engine.}
\label{tab:real-world-rotation}
\end{table}

An ordinary proposal-chain sketch is
\[
 L_r:B_k\ \longrightarrow\ \text{votes collected by }L_{r+1}
 \ \longrightarrow\ L_{r+1}:(B_{k+1},Q_r(B_k)).
\]
Neither the arrow nor the QC asserts that $B_k$ is already final. Replacing
the current engine with a chained protocol requires adopting its compatible
locking, commit and timeout rules together. Jolteon's two-chain construction
and its recovery predicates are developed as one protocol~\cite{jolteon}.

\subsection{Candidate A: complete a fixed boundary before normal handoff}
\label{sec:local-boundary-design}
Introduce a local segment index $e$ and a per-segment application quota
$K_{\rm loc}\geq1$, distinct from the cross-region quota $K$. A segment may
contain several recovery views $v$. If its first application height is
$b_e$ and the current term's last allowed height is $b_{\rm term}$, define
\[
 \bar b_e=\min\{b_e+K_{\rm loc}-1,b_{\rm term}\},\qquad
 q_e=\bar b_e-b_e+1.
\]
These quantities are derived from the authenticated opening state. Every
recovery view of segment $e$ retains the same endpoint $\bar b_e$.
Normal completion moves to $(e+1,0)$; a failure moves to $(e,v+1)$.
For example, the scheduled leader could be
\[
 \ell_i(e,v)=\mathcal L_i[(e+v)\bmod n_i].
\]
This is a candidate scheduling rule. After recovery, the next segment's
scheduled leader can coincide with the last recovery leader; a requirement
to exclude that identity would need a different authenticated schedule.

\begin{mdframed}[linewidth=.6pt,skipabove=9pt,skipbelow=9pt]
\small\noindent\textbf{Design sketch A: final-boundary local rotation.}
\begin{enumerate}[label=\textbf{A\arabic*.}]
\item \textbf{Bound admission.} Validate the complete proposed ancestry and
reject an APP beyond $\bar b_e$ in this segment, including attempts in a later
recovery view. The local counter does not reset on failure.
\item \textbf{Finish.} After the terminal APP is proposed, preserve its
commit timer. Fresh descendants may be DRAIN blocks, following the existing
completion adapter. Reaching the quota emits no BLAME and cancels no timer.
\item \textbf{Recover if necessary.} A genuine progress failure or
equivocation invokes the existing recovery path within segment $e$.
If NEWVIEW inherits the terminal APP without a running timer, a fresh DRAIN
descendant can complete it. The remaining application quota stays zero.
\item \textbf{Export a final boundary.} Any holder may disseminate the
terminal receipt $F_{\bar b_e}$, anchor $A_{\bar b_e}$ and missing data.
Completion does not require a release signature from the old leader.
\item \textbf{Open the successor segment.} Provided the region is still
authorized and the cross-region term has not ended, use the verified anchor
as the next segment's initial checkpoint. Conservatively retain the existing
NEWVIEW initialization. The first fresh APP has height $\bar b_e+1$.
\end{enumerate}
\end{mdframed}

The conservative realization treats each segment as a separately named
internal SMR instance, for example with instance identifier
$s_e=\Hash(\Tag{SEGMENT},s,e,A_{b_e-1})$. Its signing and timer contexts
include this identifier. A normal opening initializes the next instance
from a final checkpoint; it is not a timer-free direct commit inside the
old instance. Global application and mini-block heights, regional authority,
the outer term counter and its remaining quota continue unchanged.

This separation matters for DRAIN. The old physical log may extend beyond
the terminal application block through control-only descendants. Those
records are retained in the old instance; the new instance's root represents
the finalized application checkpoint. The proposed composition promises
continuity of the application projection. It does not silently discard
committed physical entries from a single continuing SMR log. A design that
instead retains one physical log must specify and certify its physical
terminal anchor as well.

The hard endpoint and the final receipt do different work. The endpoint
excludes every old-segment APP beyond $\bar b_e$; the receipt fixes the
contents through that endpoint. A receipt for a merely local ``latest
block'', without a common endpoint, supplies no closure guarantee.
Normal closure, delayed old-view recovery, and acceptance of the new
instance still require the composition argument listed below.

\subsection{Candidate A+: replace initialization with an explicit opening certificate}
\label{sec:ready-open-design}
Once the terminal anchor is final, the original leader need not release the
segment and a new leader need not discover an unknown highest old-view
anchor. An optional optimization is to make installation of the common
checkpoint explicit. Define the candidate boundary statement and evidence
\[
 \begin{aligned}
 R_e&=(\mathit{cfg},s,e,e+1,\bar b_e,A_{\bar b_e},\ell_i(e+1,0)),\\
 \mathcal H_e&=(R_e,F_{\bar b_e},\text{required ancestry and authority evidence}).
 \end{aligned}
\]
The hashed statement uses canonical block identities, not the encoding of
a particular signature list. Thus two encodings of the same final receipt
cannot create competing checkpoint identities.

After validating $\mathcal H_e$ and acquiring the data, a member may persist
the next instance's root and sign
\[
 \Sig_S(\Tag{READY},s,e+1,\Hash(R_e)).
\]
READY messages are disseminated so a failed collector cannot permanently
withhold the opening evidence. A candidate opening certificate $O_{e+1}$
contains matching READY signatures of weight $>2W_i/3$. Signers endorse at
most one root for this next-instance identifier, with that state durable
before sending. The next leader's first proposal carries $O_{e+1}$ and
extends exactly the installed anchor.

This is a \emph{new initialization rule}, not an ordinary old-view vote QC.
The first-proposal predicate must explicitly accept the opening certificate;
subsequent proposals return to the chosen engine's normal parent-QC rule.
The existing $>W_i/3$ final receipt establishes past commitment, whereas
the proposed $>2W_i/3$ READY set attests to checkpoint installation. Neither
threshold can be substituted for the other merely because both objects are
called certificates. READY also cannot prove closure of an unrestricted old
stream; the fixed boundary and its finality are prerequisites.

Before any opening certificate exists, missing terminal evidence keeps the
old segment in completion/recovery. After a valid opening is disseminated,
failure of its scheduled leader must enter recovery in the new instance at
the same installed root. Entry skew, timeout start points, certificate
availability, and the race between these cases remain explicit design
obligations. Merely collecting READY signatures is not a complete pacemaker.

\begin{figure}[htbp]\centering
% Discussion sketches; causal ordering only, not a time scale.
\begin{tikzpicture}[
  x=1cm,y=1cm,font=\footnotesize,
  box/.style={draw,rounded corners=2pt,align=center,text width=2.65cm,
              minimum height=.93cm,inner sep=4pt},
  flow/.style={-{Stealth[length=1.8mm]},semithick}
]
\node[anchor=west,font=\small\bfseries] at (-1.48,2.03)
  {A/A+: finish the fixed boundary, then open the next segment};
\node[box] (a) at (0,1.14) {Terminal APP\\proposal / QC};
\node[box] (b) at (3.45,1.14) {Complete timer;\\final receipt};
\node[box] (c) at (6.9,1.14) {Initialize, or\\READY / OPEN};
\node[box] (d) at (10.35,1.14) {Next leader's\\first fresh APP};
\draw[flow] (a)--(b); \draw[flow] (b)--(c); \draw[flow] (c)--(d);
\node[anchor=west,align=left] at (-1.48,.23)
  {Quota ends new APP admission; it does not cancel the terminal commit timer.};

\node[anchor=west,font=\small\bfseries] at (-1.48,-.65)
  {B: adopt an engine with a pipeline spanning leaders};
\node[box] (e) at (0,-1.55) {$L_r$ proposes\\$B_k$};
\node[box] (f) at (3.45,-1.55) {$L_{r+1}$ collects\\votes / parent QC};
\node[box] (g) at (6.9,-1.55) {$L_{r+1}$ proposes\\$B_{k+1}$};
\node[box] (h) at (10.35,-1.55) {Later evidence\\finalizes ancestors};
\draw[flow] (e)--(f); \draw[flow] (f)--(g); \draw[flow] (g)--(h);
\node[anchor=west,align=left] at (-1.48,-2.46)
  {QC advances leadership; the selected engine's distinct commit rule establishes finality.};
\end{tikzpicture}

\caption{Two optimization levels. Candidate A/A+ waits for the fixed terminal
application receipt; A+ changes only the subsequent opening procedure.
Candidate B permits normal leader advancement before finality by adopting
a compatible consensus engine. Arrows show causal dependencies, not timing
to scale. These are discussion sketches, not additional established protocols.}
\label{fig:local-rotation-options}
\end{figure}

\subsection{Candidate B: a pipeline whose normal progression crosses leaders}
\label{sec:qc-rotation-design}
If waiting for terminal finality at every local boundary is too expensive,
the more substantial optimization is a regional engine with QC-driven
normal leader progression, as in the implementations above. A certified
parent can then enable the next scheduled leader while an earlier block is
still progressing toward the engine's finality condition. The engine's
final output, rather than a bare vote QC, must feed the existing stream
receipt and proposer interfaces.

This cannot be implemented in the present synchronous engine by merely
changing the leader after each QC. Step S4 requires a timer to expire while
the member is still in its view, and S1 requires a current-view parent QC.
The preservation proof also uses quit propagation and the wait before the
next entry lock. A direct cross-leader QC transition changes all three
dependencies. Carrying an old timer across such a transition without a new
proof is equally unsupported.

There are two possible research routes. One is to adopt a particular
partially synchronous chained HotStuff-family engine and prove that its
finality/export interface implements the regional service. The other is
to keep synchronous finality while adopting a purpose-built rotating-leader
construction. Abraham, Nayak and Shrestha give a synchronous construction
with responsive proposal progression across leaders and $2\Delta+O(\delta)$
good-case commitment latency~\cite{rotatingsync}. Its epoch synchronization,
certificate handling and commit timers are part of that construction.
Its count-based assumptions and its lower bound for
$n/3\leq f<n/2$ do not directly establish a lower bound or an upper bound
for this manuscript's stake-weighted, below-one-third model.

Next-leader vote collection can reduce ordinary message fanout, but makes
QC availability at a faulty successor a design concern. Backup collection
and reconstruction from retained votes are relevant optimizations to study.
The present all-to-all regional votes already disseminate this information;
removing that dissemination changes both communication and recovery.
Similarly, aggregate signatures can reduce repeated certificate payloads,
but signer identification, voting-weight checks, setup and retrieval costs
must remain in the analysis. None of these changes is included in the
baseline costs of \cref{sec:complexity}.

\subsection{Timing and communication to measure for the candidates}
Let $T_{i,e}^{\rm boundary}$ run from the terminal APP proposal until its
final receipt is assembled. It equals the baseline $T_i^{\rm receipt}$ on
the healthy direct-commit path, with additional completion/recovery costs
otherwise. Let $T_{i,e}^{\rm local\text{-}xfer}$ include remaining receipt/data
dissemination and readiness at the same regional membership. It must not
be set to zero just because the membership is unchanged. For Candidate A,
the terminal-to-next-proposal interval is
\[
 H_{i,e}^{A}=T_{i,e}^{\rm boundary}
            +T_{i,e}^{\rm local\text{-}xfer}+T_{i,e+1}^{\rm init},
 \qquad T_{i,e+1}^{\rm init}\leq2\Delta_R+2\delta_R
\]
under the baseline honest-leader initialization and readiness premises.
For Candidate A+, define $T_{i,e}^{\rm open}$ from the same assembled-receipt
endpoint through dissemination, READY collection, opening and the first
new proposal. Then
\[
 H_{i,e}^{A+}=T_{i,e}^{\rm boundary}+T_{i,e}^{\rm open}.
\]
With cached data, an honest successor and a fully specified opening
pacemaker, $T_{i,e}^{\rm open}=O(\delta_R)$ is a design target. It is not a
proved bound in this version, and it does not remove the preceding
$2\Delta_R$ terminal commitment wait.

For homogeneous fault-free segments of $q$ APP batches, with no outer-term
boundary in the measurement window, the same proposal-gap accounting gives
\[
 \bar I_{\rm loc}=\frac{(q-1)I_i+H_{\rm loc}}{q},\qquad
 \lambda_{\rm loc}=\frac{q}{(q-1)I_i+H_{\rm loc}}.
\]
For a finite outer term, sum its actual internal, local-boundary and
cross-region gaps. A coincident local and cross-region endpoint has one
successor opening, not two sequential handoffs. Failed attempts and all
DRAIN blocks retain their actual costs.

Candidate A pays the already analyzed initialization exchanges for every
local segment. In A+, all-to-all READY dissemination takes $O(n_i^2)$
messages. With cached boundary evidence, its short signatures alone cost
$O(n_i^2\kappa)$ bits; repeatedly attaching or forwarding explicit
$O(n_i\kappa)$-bit certificates can restore an $O(n_i^3\kappa)$ control
term. Missing data is additional. Therefore a shorter opening delay does
not automatically imply linear communication. Candidate B requires fresh
accounting from the selected engine, including its failure path and the
post-finality receipt exports used here.

\subsection{Proof obligations and evaluation priorities}
\label{sec:rotation-obligations}
Before any candidate becomes normative, it must discharge the following
obligations rather than inherit them from its name or from an implementation
with different assumptions.
\begin{enumerate}[label=\textbf{P\arabic*.}]
\item \textbf{Unique boundary and successor.} Every accepted terminal
receipt for segment $e$ determines the same application anchor; normal
opening and delayed recovery cannot authorize competing successors or
application positions. The claim includes the physical-log/application
projection distinction.
\item \textbf{Preservation under overlap.} A node learning a terminal
receipt while peers are recovering cannot cause loss of completed honest
work. READY persistence, late votes, timer retirement and data retention
must be reconciled explicitly. Old release signatures are unnecessary.
\item \textbf{Completion before normal advancement.} For A/A+, quota
completion keeps pending timers alive and recovery does not repeatedly
reset the segment. Establish eventual terminal finality and opening even
if the scheduled successor or a collector fails. For B, establish the
corresponding chain-based progress property instead.
\item \textbf{Leader schedule compatible with weights.} Faulty stake below
$W_i/3$ does not imply that fewer than one third of identities are faulty.
For example, one honest identity may hold $0.8W_i$ and three faulty identities
the remaining $0.2W_i$. A once-per-identity cycle then has no consecutive
honest leaders. A chained engine whose progress argument needs such a run
cannot infer it from the current assumptions. Its schedule, leader tenure
or probabilistic progress premise must supply what its proof requires.
\item \textbf{Outer authority and consumer compatibility.} A local opening
does not reset $K$, bypass an effective L1 retirement or alter delivered
prefixes. The rollup verifier must authenticate local-segment succession if
it becomes part of signing authority. Execution, DA and L1 settlement keep
their separate meanings.
\item \textbf{Measured endpoints.} Measure proposal spacing, terminal
finality, first successor receipt, recovery delay and bytes separately.
Include $K_{\rm loc}=1$, small $\delta_R/\Delta_R$, missing ancestry, a
withholding QC collector, overlapping receipt and timeout delivery, and
coincident local/outer/L1 boundaries. A bounded simulation can expose bad
traces but does not replace the weighted safety and progress proofs.
\end{enumerate}

The immediate extension to formalize is Candidate A: fixed application
segments, final-boundary completion and ordinary initialization, with
BLAME reserved for recovery. A+ isolates a subsequent startup optimization.
Candidate B addresses the stronger goal of sustaining a pipeline across
leaders without draining it at every normal boundary; it is an engine
change and should receive its own specification and proof before replacing
the regional baseline.
